\scanpage{7}
\mspage{5}

éventuellement muni de la sous-catégorie pleine de tous les systêmes loc
sur X, et du foncteur cohomologique sur $D^{b}_{lc}(X)$ à valeurs dqns laddt\unsure{ë}
catégorie, et bien sûr du produit tensoriel (mais alors on sort de $D^{b}$
pour entrer dans $D^{-}$ , redactor demerdetur).

II \quad $n$-catégories, catégories $n$-uples, et (Gr)-catégories.

Appelons bi-catégorie une catégorie dans(Cat), i.e. un triple
($C^{\deleted{o}}_{\handwritten{0}}$,$C^{\deleted{1}}_{\handwritten{1}}$, \deleted{\illegible{2}o\unsure{ï|i}}$p_{1}$,$p_{2}$, $\handwritten{\pi}$ ), où $C^{o}$ est une catégorie, $p_{1}$,$p_{2}$: $C^{\deleted{1}}_{\handwritten{1}} \handwritten{\rightrightarrows} C^{\deleted{o}}_{\handwritten{0}}$ deux
foncteurs, $\handwritten{\pi}$:$(C^{\deleted{1}}_{\handwritten{1}},p_{1})\mathrm{x}_{C^{o}}(C^{\deleted{1}}_{\handwritten{1}},p_{2}) \handwritten{\longrightarrow} C^{\deleted{1}}_{\handwritten{1}}$ un foncteur, avec les axiomes
habituels. Interprétant (Cat) comme la sous-catégorie pleine de (Ssimpl
des objets qui .... , et de même pour les catégories à valeurs dans une
catégorie telle (Cat), on trouve que la catégorie \inserted{(Bicat)} des (Cat)-catégories
est équivalente à la sous-catégorie pleine de la catégorie des ensembles
bi-semisimpliciaux

\[
\begin{array}{cccc}
X^{\deleted{00}}_{\handwritten{00}} & X^{\deleted{01}}_{\handwritten{01}} & \ldots & X^{\deleted{0n}}_{\handwritten{0n}}\\
X^{\deleted{10}}_{\handwritten{10}} & X^{\deleted{11}}_{\handwritten{11}} & \ldots & X^{\deleted{1n}}_{\handwritten{1n}}\\
\multicolumn{2}{l}{.\,.\,.\,.\,.\,.\,.} & &
\end{array}
\]

formée des objets qui ligne par ligne et colonne par colonne satisfoht la
condition d'exctitude familière (transformer sommes amalgamées en produi\unsure{t}
fibrés). De même la catégorie \deleted{des} (Tricat) des (Bicat)-catégories
\unsure{\deleted{'}s|\deleted{s}'}\unsure{\deleted{s}'|\deleted{'}s}interprète en t ermes d'objets tri-\begin{tabular}[b]{@{}l@{}}\inserted{semi}\\[-.6ex]\begin{tikzpicture}[x=.65cm,y=.4cm,baseline=0pt]\draw[line width=.5pt] (0,-.35)..controls(.25,.2)and(.5,.45)..(.8,.5)--(1.1,.53);\draw[line width=.5pt] (1.22,.57)..controls(1.55,.75)and(2.1,.7)..(2.3,.95)..controls(2.38,1.04)and(2.35,1.16)..(2.3,1.25);\end{tikzpicture}\end{tabular}simpliciaux, etc.

Exemple: soit C une catégorie, posons $C^{\deleted{o}}_{\handwritten{0}}$=C, $C^{\deleted{1}}_{\handwritten{1}}$=$\underline{Fl}$(C), avec structure
$p_{1}$,$p_{2}$, évidentes (d'où une (Cat)-catétorie dont les composantes semi-si\unsure{m}pliciales supérieures sont les catégories $\underline{Fl}^{n}$(C) des $n$-flêches de C ...)
L'ensemble bi-semisimplicial correspondant a pour composantes les ensembl
Hom($\handwritten{\Delta}^{n}\handwritten{\unsure{\times|\alpha}}\handwritten{\Delta}^{m}$,C). Par exemple, les objets de $X^{\deleted{11}}_{\handwritten{11}}$ sont les carrés commutatif
de flêches de C, qu'on peut interpréter(de deux manières différentes) co\unsure{m}

\margin{6}
